\subsection{Limitations}
\label{app:sec:limitations}
\noindent \textbf{Benchmark Size.}
\livecodebench{} code generation scenario currently hosts over $500$ instances from problems released since May $2023$.
To account for contamination in \deepseek{}, we only perform evaluations on problems released after the model cutoff date.
% However, for newer models, we perform evaluations on post-cutoff problems in the dataset.
This leads to only $349$ problems used in our final evaluations which might add noise due to problem set samples.
We currently estimate $1-2\%$ performance variations in \lcb{} code generation due to this issue (measured by bootstrapping $349$ sized problem sets from the full dataset).
Other scenarios, i.e. self-repair, code execution, and test output prediction comprise $238$, $188$, and $254$ problems would have similar performance variations. 
We thus recommend exercising proper judgement when comparing models with small performance differences.
Note that \humaneval{} has $164$ problems and would also struggle with similar issues.

This issue is also exacerbated for newer models, with more recent cutoff dates, as they might only have access to a smaller evaluation set.
We propose two solutions addressing this issue as we evolve \livecodebench{}.
First, we will use other competition platforms for problem collection, allowing larger number of recent problems to be added to the benchmark.
In addition, we also hope supplement this with an unreleased private test set constructed specifically for model evaluation. 
These problems will use a similar flavor to current problems and will be used when models are submitted for evaluation to the \livecodebench{} platform.
This would reduce the reliance on public accessible problems and provide a more robust evaluation of the models while providing community public access to similar problems, similar to strategies employed by popular platforms like \smalltextsc{Kaggle}.
%hidden set which can be used for all 
%First, depending on the models used, we might be able to adapt some of the older problems for evaluation.
%Second, going forward, we will collect problems from other competition platforms thus improving the counts of \textit{new} problems. Further, we can also gather a private test set which will not be released allowing robust evaluations with larger problem sizes.  

\vspace{1pt}
\noindent \textbf{Focus on \python{}.} 
\livecodebench{} currently only focuses on \python{} which might not provide enough signal about model capabilities in other languages. 
However, since we collected problem statements and serialized tests, adding new programming languages would be straightforward once appropriate evaluation engines are used.

\vspace{1pt}
\noindent \textbf{Robustness to Prompts.}
Recent works have identified huge performance variances that can be caused due to insufficient prompt.
Here, we either do not tune prompts across models or make minor adjustments based on the system prompts and delimiter tokens.
This can lead to performance variance in our results. 
Our findings and model comparison orders generalize across \livecodebench{} scenarios 
and mostly match the performance trends observed on \humaneval{} making this a less prominient issue.

This issue can be particularly observed open models on the code execution scenario with \chainot{} prompting.
Interestingly, often the open models perform even worse in comparsion to the direct code execution baseline.
Note that we used same prompts for the closed models all of which show noticable improvement from \chainot{}.
%However, we do find that open models perform poorly on the code execution task with \chainot{}, even performing worse than non-\chainot{} setups.
%Note that common prompts were used across both open and closed models during these experiments.
While the used prompts might be sub-optimal, this highlights how open-models perform worse against the closed models at performing chain-of-thought.
%This means prompts might be sub-optimal for this setting but at the same time also highlights that open models seem to be worse than closed models at performing chain-of-thought.

\vspace{1pt}
\noindent \textbf{Problem Domain.}
Programming is a vast domain and occurs in various forms such as programming puzzles, competition programming, and real-world software development.
Different domains might have individual requirements, constraints, challenges, and difficulty levels.
\livecodebench{} currently focuses on competition problems sourced from three platforms.
This might not be representative of the ``most general'' notion of \llm{} programming capabilities. 
Particularly, real-world usage of \llms{} is drawn upon open-ended and unconstrained problems rasied by users.
We therefore recommend using \livecodebench{} as a starting point for evaluating \llms{} and 
further using domain-specific evaluations to measure and compare \llms{} in specific settings as required.